\documentclass[11pt]{article}
\usepackage[margin=1in]{geometry}
\usepackage{amsmath,amssymb,amsthm,mathtools}
\usepackage[hidelinks]{hyperref}

\newtheorem{theorem}{Theorem}[section]
\newtheorem{lemma}[theorem]{Lemma}
\newtheorem{proposition}[theorem]{Proposition}
\theoremstyle{remark}
\newtheorem{remark}[theorem]{Remark}
\newcommand{\Arg}{\operatorname{Arg}}
\newcommand{\tr}{\operatorname{tr}}
\newcommand{\one}{\mathbf 1}

\title{Positive Square Energy of Every Tetracyclic Cactus}
\author{}
\date{26 July 2026}

\begin{document}
\maketitle

\begin{abstract}
We prove that every connected tetracyclic cactus $G$ of order $n$ satisfies
$s^+(G)>n$.  The sharp block-additive doubly nonnegative estimate reduces
the theorem exactly to the cycle multisets
$\{3,3,3,q\}$, with $q$ odd, and $\{3,3,5,5\}$.  For the first family we use
maximum-cycle-packing territories when $q\equiv3\pmod4$ and a shared-cut
cluster territory lemma when $q\equiv1\pmod4$.  For the second family,
ordinary bridge cuts handle disconnected shared-cut clusters; one connected
cluster has twenty exact core types.  Nineteen have induced-packet
certificates, while the common-root bouquet is handled by an exact two-state
lobe phase theorem.  All connector trees and all trees attached at core
vertices are arbitrary.  No cycle-closure monotonicity is assumed, and no
raw, coefficientwise phase certificate is asserted where it is false.
\end{abstract}

\section{Preliminaries}

All graphs are finite, simple, and undirected.  A cactus is a connected graph
whose blocks are edges or cycles.  For the adjacency eigenvalues
$\lambda_1,\ldots,\lambda_n$, write
\[
 s^+(G)=\sum_{\lambda_j>0}\lambda_j^2,
 \qquad s^-(G)=\sum_{\lambda_j<0}\lambda_j^2,
 \qquad D(G)=s^+(G)-s^-(G).
\]
The quantities were introduced in connection with spectral coloring bounds
in~\cite{AbiadEtAl}.  Liu--Tang--Zhang proved the general bound
$\min\{s^+,s^-\}\geq n-1$~\cite{LTZ}.  Akbari--Kumar--Mohar--Pragada--Zhang
conjectured the stronger conclusion $s^+(G)\geq n$ for every connected graph
with at least $n+1$ edges~\cite{AKMPZ}.  We prove that conclusion here for
tetracyclic cacti.

A connected tetracyclic graph has $m=n+3$, and therefore
\begin{equation}\label{eq:average}
 s^+(G)+s^-(G)=2m=2n+6,
 \qquad s^+(G)=n+3+\frac12D(G).
\end{equation}
For every vertex partition $(V_j)$ we use induced-subgraph superadditivity
\begin{equation}\label{eq:super}
 s^+(G)\geq\sum_j s^+(G[V_j]),
\end{equation}
proved in~\cite{AKMP}.  Cross edges are ignored in \eqref{eq:super}; no claim
that adding an edge or closing a cycle increases $s^+$ is made or needed.

For positive activities $a=(a_v)$ put
\[
 Z_F(a)=\sum_{M\text{ matching of }F}\prod_{v\text{ unmatched by }M}a_v,
 \qquad Z_F(t)=Z_F((t)_v),
\]
and normalize the characteristic polynomial by
\[
 \Psi_G(t)=i^{-n}\det(itI-A(G))=\prod_j(t+i\lambda_j).
\]
Sachs' expansion, grouped by collections $\mathcal C$ of pairwise
vertex-disjoint cycles, is
\begin{equation}\label{eq:sachs}
 \Psi_G(t)=\sum_{\mathcal C}
 \prod_{C\in\mathcal C}(-2i^{-|C|})Z_{G-V(\mathcal C)}(t).
\end{equation}
Thus cycles of lengths $1$ and $3$ modulo four have multipliers $2i$ and
$-2i$, respectively.  If $\Theta_G(t)$ is the continuous argument of
$\Psi_G(t)$ tending to zero as $t\to\infty$, then
\begin{equation}\label{eq:coulson}
 D(G)=-\frac4\pi\int_0^\infty t\Theta_G(t)\,dt.
\end{equation}
This follows by integrating
$\lambda|\lambda|=(2/\pi)\int_0^\infty
\lambda^3/(\lambda^2+t^2)\,dt$, summing, using $\sum_j\lambda_j=0$, and
integrating by parts.

Arbitrary attached trees can be eliminated exactly.  Orient a rooted branch
toward its core attachment.  If $T_{u\to p}$ is the subtree below $u$, then
\begin{equation}\label{eq:BP}
 Q_{u\to p}=\frac{Z_{T_{u\to p}}(t)}{Z_{T_{u\to p}-u}(t)}
 =t+\sum_{w\text{ child of }u}Q_{w\to u}^{-1}\geq t.
\end{equation}
After elimination, every retained core vertex has an activity
$a_v=t+\sum Q_{u\to v}^{-1}\geq t$, and every Sachs term has the same positive
factor.  Splitting a matching according as its root is unmatched or matched
to one child proves \eqref{eq:BP}, including the terms in which a selected
Sachs cycle deletes the attachment vertex.

We use the following lower-rank results.  They are proved for arbitrary
connectors and arbitrary attached trees in the cited companion manuscripts;
the odd-unicyclic sign theorem originates with Ning--Zeng~\cite{NingZeng}.
Put
\[
 \delta_q=\sec(\pi/q)-1\quad(q\equiv1\pmod4),
 \qquad \delta_5=\sqrt5-2.
\]

\begin{lemma}[Packet bounds]\label{lem:packets}
Let $H$ be a connected cactus of order $h$.
\begin{enumerate}
\item If every cycle is $3\pmod4$ and the cycle-packing number is at most two,
then $D(H)>0$.  In particular, a packet with $r=1$ or $2$ triangles has
$s^+(H)>h+r-1$.
\item A unicyclic $C_q$ packet, $q\equiv1\pmod4$, has
$s^+(H)\geq h-\delta_q$.
\item A $C_3,C_q$ bicyclic packet, $q\equiv1\pmod4$, has
$s^+(H)>h+1-\delta_q$.
\item Every bicyclic cactus has $s^+(H)\geq h$.
\item A $C_3,C_3,C_q$ tricyclic packet with $q$ odd has $s^+(H)>h$.
If the two triangles share a cut vertex and $q=5$, the shared-triangle phase
calculation supplies the stronger bound $s^+(H)>h+2-\delta_5$.
\item A $C_3,C_5,C_5$ tricyclic packet has $s^+(H)>h$.
\end{enumerate}
Every assertion remains valid with arbitrary trees attached anywhere.
\end{lemma}

\begin{proof}
For (1), \eqref{eq:sachs} has strictly negative imaginary part: singleton
cycle terms are negative imaginary, and, because at most two cycles can be
selected, all remaining terms are real.  Equations
\eqref{eq:coulson}--\eqref{eq:average}, in the appropriate cyclomatic rank,
give the claim.  Tree elimination gives a unicyclic $C_q$ packet the phase
of $Z_{C_q}(a)+2i$.  Since $a_v\geq t$, its phase is at most that of the bare
cycle, for which
\[
 D(C_q)=-2\delta_q,
 \qquad s^+(C_q)=q+1-\sec(\pi/q).
\]
This proves (2).  The mixed $C_3,C_q$ product-subpartition comparison proves
(3), while (4) is the complete bicyclic-cactus theorem
\cite{AllBicyclic}.  The complete tricyclic-cactus theorem
\cite{AllTricyclic} gives (5)--(6); its shared-triangle $\{3,3,5\}$ phase
calculation gives the displayed stronger bound under the stated incidence
hypothesis.  These arguments use actual induced packets or exact matching
phases, never cycle-closure monotonicity.
\end{proof}

For later bookkeeping, Table~\ref{tab:credits} gives the credits
$s^+(H)-h$ used in the finite certificates.  A word ``strict'' means that
equality with the displayed lower bound is excluded.
\begin{table}[ht]
\centering\small
\begin{tabular}{lll}
cycle content & credit & status\\\hline
tree &$-1$& weak\\
$3$ &$0$& strict\\
$3,3$ &$1$& strict\\
$5$ &$2-\sqrt5=-\delta_5$& weak\\
$3,5$ &$3-\sqrt5=1-\delta_5$& strict\\
$5,5$ &$0$& weak\\
$3,3,5$ with shared triangles &$4-\sqrt5=2-\delta_5$& strict\\
$3,5,5$ &$0$& strict
\end{tabular}
\caption{Proved packet credits, all allowing arbitrary attached trees.}
\label{tab:credits}
\end{table}

\section{The exact DNN residual classification}

For a connected graph with an edge define the Liu--Tang--Zhang functional
\[
 \kappa(G)=\sup_{\substack{M\succeq0,\ M\geq0\\\one^TM\one>0}}
 \frac{4(\sum_{uv\in E(G)}\sqrt{M_{uv}})^2}{\one^TM\one}.
\]

\begin{lemma}[Sharp cactus DNN bound]\label{lem:dnn}
Suppose a cactus has $b$ bridge blocks.  If its cyclic blocks are
$C_{\ell_1},\ldots,C_{\ell_r}$, then
\[
 \kappa(G)=b+\sum_{j=1}^r\kappa(C_{\ell_j}),\qquad
 \kappa(C_\ell)=\begin{cases}
 \ell,&\ell\text{ even},\\
 \ell\sec^2(\pi/(2\ell)),&\ell\text{ odd},
 \end{cases}
\]
and $s^-(G)\leq\kappa(G)$.
\end{lemma}

\begin{proof}
We recall the sharp argument from~\cite{SharpDNN}.  Folding a nonedge entry
$M_{uv}$ into the diagonal by adding
$M_{uv}(e_u-e_v)(e_u-e_v)^T$ preserves the denominator and edge entries.
On a cycle, rotational averaging and concavity of the square root reduce the
problem to common edge entry $x$ and row sum $r$.  A Fourier eigenvalue gives
$r\geq4x$ for even $\ell$ and
$r\geq2(1+\cos(\pi/\ell))x$ for odd $\ell$.  Equality is attained by
$x(2cI+A(C_\ell))$, with $c=1$ or $c=\cos(\pi/\ell)$.

At a one-vertex sum, fold first and use Gram vectors.  The spans on opposite
sides are orthogonal.  Splitting the cut vector into its two span projections
splits both the edge numerator and $\one^TM\one$; Cauchy--Schwarz gives
subadditivity.  Orthogonal gluing of near-optimizers, scaled so
$T_1/\kappa_1=T_2/\kappa_2$, gives the reverse inequality.  Hence $\kappa$ is
block additive and $\kappa(K_2)=1$.

Finally write $A=P-B$ for the positive and negative spectral parts.  The
Schur product theorem makes $B\circ B$ doubly nonnegative, while
$s^-/2=-\sum_{uv\in E}B_{uv}\leq\sum_{uv\in E}|B_{uv}|$ and
$\one^T(B\circ B)\one=s^-$.  The definition of $\kappa$ gives
$(s^-)^2\leq\kappa(G)s^-$.
\end{proof}

Put $\epsilon_\ell=0$ for even $\ell$ and
\[
 \epsilon_\ell=\ell\tan^2\!\left(\frac\pi{2\ell}\right)
 \quad(\ell\text{ odd}).
\]
Then $\kappa(C_\ell)=\ell+\epsilon_\ell$,
\begin{equation}\label{eq:eps}
 \epsilon_3=1,
 \qquad \epsilon_5=5-2\sqrt5,
 \qquad \epsilon_5+\epsilon_7<1,
\end{equation}
and $\epsilon_\ell$ strictly decreases through odd integers.  Indeed, after
$u=\pi/(2x)$ the monotonicity is that of $\tan^2u/u$, whose derivative is
positive because $2u>\sin u\cos u$.  Also
$\cos(\pi/7)>7/8$ gives
$\epsilon_7<7/15<2\sqrt5-4=1-\epsilon_5$.

For four cycle lengths $\ell_1,\ldots,\ell_4$, block counting gives
$b+\sum_j\ell_j=n+3$.  Lemma~\ref{lem:dnn} and the energy identity yield
\begin{equation}\label{eq:dnn-bound}
 s^+(G)\geq n+3-\sum_{j=1}^4\epsilon_{\ell_j}.
\end{equation}

\begin{proposition}[Exact residual list]\label{prop:residual}
The sharp DNN estimate \eqref{eq:dnn-bound} proves $s^+(G)\geq n$ unless the
cycle multiset is
\begin{equation}\label{eq:residual}
 \{3,3,3,q\}\quad(q\geq3\text{ odd}),
 \qquad\text{or}\qquad \{3,3,5,5\}.
\end{equation}
These and only these multisets have $\sum_j\epsilon_{\ell_j}>3$.
\end{proposition}

\begin{proof}
Four triangles give the first family with $q=3$.  With exactly three
triangles, the excess sum exceeds three precisely when the fourth length is
odd, giving all remaining $\{3,3,3,q\}$.  With exactly two triangles, the two
other lengths must be odd and their excesses must sum to more than one.
Equation \eqref{eq:eps} and strict monotonicity leave exactly $5,5$.  With at
most one triangle, the sum is at most
$1+3\epsilon_5<3$.  Even lengths contribute zero.  This exhausts all
unordered length multisets.
\end{proof}

The list is a classification of what the sharp DNN relaxation leaves, not a
list of counterexamples and not a claim that the spectral inequality in the
DNN proof is sharp.

The only nonresidual case in which \eqref{eq:dnn-bound} is not already strict
has three triangles and one even cycle.  We record the strict argument needed
for that boundary case.

\begin{proposition}[Three triangles and one even cycle]\label{prop:333even}
If the cyclic blocks of $G$ are $C_3,C_3,C_3,C_q$ with $q$ even, then
$s^+(G)>n$.
\end{proposition}

\begin{proof}
Suppose first that the shared-cut graph is disconnected.  In the reduced
cluster tree choose a leaf cluster not containing $C_q$ and cut its first
actual bridge toward the rest.  This leaf is a connected shared-cut cluster of
$r\in\{1,2,3\}$ triangles.  By Lemma~\ref{lem:packets}(1), its surplus is
strictly greater than $r-1$.  Its connected complement has rank $4-r$ and is
nonnegative by the unicyclic even-cycle bound when $r=3$ and by
Lemma~\ref{lem:packets}(4) or the complete tricyclic theorem when $r=2$ or
$r=1$.  Induced-subgraph superadditivity therefore gives strict positive total
surplus.

Now suppose all four cycles form one shared-cut cluster, and inspect the
$C_q$ node in its cycle--cut incidence tree.  If it is a leaf, open a vertex
of $C_q$ private with respect to cyclic blocks.  The opened territory is a
nonempty tree of surplus $-1$, while deleting the leaf cycle node leaves the
three triangles in one connected shared-cut cluster.  Their packing number is
at most two, so Lemma~\ref{lem:packets}(1) gives surplus greater than $2$;
the total is greater than $1$.

If the $C_q$ node has degree $d\geq2$, split it into $d$ nonempty proper
consecutive intervals, one for each component of the incidence tree after
deleting that node.  Every component contains at least one triangle.  If their
triangle counts are $r_1,\ldots,r_d$, the resulting induced territories have
surplus strictly greater than $r_i-1$, and hence their total surplus is
strictly greater than
\[
 \sum_i(r_i-1)=3-d\geq0.
\]
All connector remnants and hanging trees follow their unique attachment, as in
Lemma~\ref{lem:333qterritory}.  Thus every incidence has strict positive
surplus.
\end{proof}

\section{Three triangles and one odd cycle}

We first record the territory statements needed here.  A cyclic territory is
an induced connected vertex set containing at least one cyclic block.  Every
component outside a chosen cyclic core has a unique core root: two roots and a
path through the component would create another cycle.  Assigning each whole
rooted component to the territory containing its root extends every core
partition over arbitrary trees.

\begin{lemma}[Maximum-packing territories]\label{lem:maxterritory}
In a cactus, let $\mathcal P$ be a maximum collection of pairwise
vertex-disjoint cyclic blocks.  There is a partition into connected induced
territories, one containing each member of $\mathcal P$, such that every
territory has cycle-packing number at most two.
\end{lemma}

\begin{proof}
Use the block-cut tree, subdividing bridge-only paths as necessary.  Delete
edges of this tree so that each resulting component contains exactly one
selected cycle node; assign a tie component consistently toward either
neighbor.  Lifting the components gives connected induced territories, and
hanging branches follow their unique attachment.  If one territory contained
three pairwise disjoint cyclic blocks, none could lie outside that territory
and each is disjoint from the selected cycles in all other territories.
Replacing its one selected block by those three would enlarge $\mathcal P$,
a contradiction.  Thus its packing number is at most two.  Notice that this
is an induced partition argument, not the assertion that closing any of the
cut cycles is spectrally monotone.
\end{proof}

For the opposite residue we need a small block-tree fact.  Join two cyclic
blocks when they share a cut vertex; components of this relation are called
shared-cut clusters.

\begin{lemma}[The $C_{333q}$ cluster territories]\label{lem:333qterritory}
Let a cactus have cyclic blocks $C_3,C_3,C_3,C_q$, where
$q\equiv1\pmod4$, and let $\delta=\sec(\pi/q)-1<1$.
Its vertices admit induced connected packet territories whose total proved
credit is positive or is a strict lower bound of zero.  Every packet used has
cycle-packing number at most two, is a $C_q$, $C_3,C_q$, or
$C_3,C_3,C_q$ packet, or contains only triangles; tree packets occur only in
the one-attachment connected case.
\end{lemma}

\begin{proof}
First suppose that the shared-cut graph is disconnected.  Contract every
shared-cut cluster and every maximal bridge-only connector; the resulting
cluster incidence is a tree.  Let $t$ be the number of triangles in the
cluster containing $C_q$.  Cutting connector edges partitions this tree into
induced connected cluster territories, with each intervening connector put
wholly on one of its two sides.  If $t=1$ or $2$, use the rank-two or
rank-three $C_q$ territory and the remaining triangular clusters.  Its credit
is respectively greater than $1-\delta$ or greater than $0$, while each
remaining triangular cluster is put in a strict triangular packet (one per
cluster).  If
$t=0$ and some other cluster contains at least two triangles, that cluster has
credit greater than at least $1$, which pays the $C_q$ singleton deficit
$-\delta$; a three-triangle cluster still has packing number at most two,
because its shared-cut graph is connected and hence two triangles meet.  All
other triangular clusters have nonnegative strict credit.  If
$t=0$ and all three triangular clusters are singletons, take the unique path
in the cluster tree from the $C_q$ cluster to a nearest triangular cluster,
put that entire connector into their $C_3,C_q$ bicyclic packet, and leave the
other two triangles in strict triangular packets.  Its credit is greater than
$1-\delta>0$.

Now suppose that all four cyclic blocks form one shared-cut cluster.  Remove
the $C_q$ block node from its block-cut tree.  Group together the triangle
blocks in each remaining component: triangles attached indirectly through
other triangles therefore belong to the same group.  The nonempty groups
attach to distinct vertices $v_1,\ldots,v_d$ of $C_q$; write $r_i$ for their
triangle counts, so $r_i\geq1$ and $\sum_i r_i=3$.

Assume first that $d\geq2$.  Fix a cyclic orientation of $C_q$.  Removing the
attachment vertices leaves exactly $d$ (possibly empty) open path components.
Assign every internal vertex of each such path, together with every branch
rooted at an internal vertex, to the group at its preceding endpoint.  Thus
the packet $H_i$ consists of the $i$th triangle component group, $v_i$, and
the interiors assigned to $v_i$.  The first edge of every nonempty assigned
arc joins its interior to $v_i$; an empty arc requires no vertices.  The edge
at the other end goes to the next packet.  Consequently the $H_i$ are
connected induced subgraphs and partition all vertices.  No packet contains
all of $C_q$, and the complete cycles of $H_i$ are exactly its $r_i$
triangles.  Since $d\geq2$ and the $r_i$ are positive with sum three, each
$r_i\leq2$, so Lemma~\ref{lem:packets}(1) gives packet credit strictly greater
than $r_i-1$.  The total credit is therefore strictly greater than
$\sum_i(r_i-1)=3-d\geq0$.

If $d=1$, put the entire triangle component group, including its common
attachment vertex $v_1$, in one packet, and put the path $C_q-v_1$ in a tree
packet.  The triangle group is connected in the shared-cut graph, so two of
its three triangles meet and its cycle-packing number is at most two.  The
first packet therefore has credit greater than $2$ by
Lemma~\ref{lem:packets}(1), while the tree packet has credit at least $-1$;
their total credit is greater than $1$.  This includes the common-root
bouquet.  In every case each off-core component is assigned whole to the
packet containing its unique root, including components rooted on assigned
arc interiors, so arbitrary attached trees preserve connectivity, inducedness,
and cycle content.

This is purely block-tree bookkeeping.  In particular, the proof does not use the
tempting raw comparison $\Phi=2R-Z_qI$: at $q=5$ that polynomial has negative
coefficients for eight of the eleven shared core types, even after substituting
$a_v=t+y_v$.  The exact failed-certificate audit is listed in
Appendix~\ref{app:reproduce}; no false coefficientwise $\Phi$ claim enters
this lemma.
\end{proof}

\begin{proposition}\label{prop:333q}
If the cyclic blocks of $G$ are $C_3,C_3,C_3,C_q$ with $q$ odd, then
$s^+(G)>n$.
\end{proposition}

\begin{proof}
If $q\equiv3\pmod4$, apply Lemma~\ref{lem:maxterritory}.  Every territory has
only $3\pmod4$ cycles and packing number at most two, so
Lemma~\ref{lem:packets}(1) gives $D>0$ on each.  A territory with $r\geq1$
cycles and order $h$ consequently has
$s^+>h-1+r\geq h$.  Superadditivity proves $s^+(G)>n$.

Let $q\equiv1\pmod4$.  Lemma~\ref{lem:333qterritory} and
Table~\ref{tab:credits} give positive total credit, or zero total credit with
at least one strict packet; either conclusion and \eqref{eq:super} give
$s^+(G)>n$.  The common-root theorem below independently gives the stronger
bound $s^+(G)\geq n+3-\delta_q$ for the bouquet.
\end{proof}

\section{Two triangles and two pentagons}

\subsection{Disconnected shared-cut clusters}

Assume first that the shared-cut graph has more than one component.  Its
contracted cluster incidence is a tree.  Deleting a connector bridge gives
vertex sets of connected induced subgraphs.  No connector length is bounded:
each bridge-tree component stays wholly on one side of the cut.

\begin{proposition}\label{prop:3355disconnected}
If the $C_3,C_3,C_5,C_5$ blocks do not form one shared-cut cluster, then
$s^+(G)>n$.
\end{proposition}

\begin{proof}
Write $\delta=\delta_5=\sqrt5-2$.  The distributions of four cycle blocks
among clusters are as follows.

For a $3+1$ distribution, isolate the singleton leaf cluster.  A singleton
triangle is strict with zero credit.  If it is a pentagon, retain the hostile
partition exactly as
\begin{equation}\label{eq:hostile}
 \{3,3,5\}\mid\{5\}.
\end{equation}
Let $H$ be the internal $\{3,3,5\}$ cluster and let the other pentagon be
joined to it by an arbitrary connector tree.  If the two triangles of $H$
meet, the shared-triangle phase bound in Lemma~\ref{lem:packets}(5) gives
credit greater than $2-\delta$ on $H$, while the external pentagon has credit
at least $-\delta$.  Thus the total credit is greater than
$2-2\delta=6-2\sqrt5>0$.

Suppose instead that the two triangles $T_1,T_2$ are disjoint.  Because $H$
is one shared-cut cluster, its pentagon $P$ meets both triangles, at distinct
vertices $x_1,x_2$.  Partition its core into
\[
 U_1=V(T_1),\qquad U_2=V(T_2)\cup\bigl(V(P)\setminus\{x_1\}\bigr).
\]
These are induced and connected: $P-x_1$ is a path containing $x_2$.
Whichever packet owns the internal root of the connector absorbs the full
connector and the external $C_5$.  On $P$, the root is owned by $U_1$ exactly
when it is $x_1$, and otherwise by $U_2$; a root on either triangle is owned
by the packet containing that triangle.  Hence the enlarged packet has cycle
content $C_3,C_5$ and credit greater than $1-\delta$, while the other triangle
packet has strict credit $0$.  The total is greater than
$1-\delta=3-\sqrt5>0$, and hence $s^+(G)>n+3-\sqrt5$.  Every branch at the
connector root or at an internal connector vertex is assigned with that whole
connector side, so no restriction on connector length or rooted branches is
used.  This is the
hostile exact construction; a uniform strong bound for arbitrary
$\{3,3,5\}$ incidence is neither asserted nor needed.

For $2+2$, the cut is either
$\{3,3\}\mid\{5,5\}$, with total credit strictly greater than $1$, or
$\{3,5\}\mid\{3,5\}$, with total credit
$2(1-\delta)=6-2\sqrt5>0$.

For $2+1+1$, cut the two appropriate connector edges.  According to the
two-cycle cluster, the credit sums are
\[
 \{3,3\}\mid\{5\}\mid\{5\}:\quad 1-2\delta=5-2\sqrt5>0,
\]
\[
 \{3,5\}\mid\{3\}\mid\{5\}:\quad 1-2\delta>0,
 \qquad
 \{5,5\}\mid\{3\}\mid\{3\}:\quad 0\ \hbox{strict}.
\]
For $1+1+1+1$, isolate a triangular leaf if one exists.  Otherwise the two
triangles are internal and a pentagon is a leaf; deleting that leaf leaves a
connected $\{3,3,5\}$ territory, so the two structural cases following
\eqref{eq:hostile} apply.  These four integer
distributions exhaust every disconnected cluster tree.
\end{proof}

\subsection{One shared-cut cluster: nineteen packet certificates}

The bare cyclic core has $13$ vertices and $16$ edges.  Recursively adjoining
leaf cycle blocks, then quotienting independently by exact graph isomorphism
and by canonical colored-cycle coding, gives exactly twenty types.  Label the
cycles $T_0,T_1,P_2,P_3$.  An incidence string such as $012+23$ records the
sets of cycles meeting at each shared cut vertex.

For each core, the exact audit tests every nonempty vertex subset.  It retains
a subset only when the induced graph is connected, its complete-cycle content
is a row of Table~\ref{tab:credits}, and its cyclomatic number equals the
number of retained listed cycles.  A subset dynamic program anchors each new
part at the least remaining vertex, thereby exhausting every set partition
without packet-order duplicates.  Credits are pairs $a+b\sqrt5$ and are
compared by integer sign and squaring tests, never floating point.

Table~\ref{tab:3355} is the complete certificate ledger.  Vertex labels refer
to the canonical cycles printed by the script in Appendix~\ref{app:reproduce}.
Semicolons separate induced packets; every unlisted packet vertex is absent
only from the abbreviated cycle-content column, not from the explicit vertex
sets.

\begin{table}[p]
\centering\scriptsize
\begin{tabular}{rlllp{2.3cm}}
type&incidence&packet vertex sets&cycle contents&credit\\\hline
2&012+03&$\{1,9,10,11,12\};\{0,2,3,4,5,6,7,8\}$&$P_3;T_1P_2$&$5-2\sqrt5$, strict\\
3&012+23&$\{5,9,10,11,12\};\{0,1,2,3,4,6,7,8\}$&$P_3;T_0T_1$&$3-\sqrt5$, strict\\
4&012+23&$\{6,9,10,11,12\};\{0,1,2,3,4,5,7,8\}$&$P_3;T_0T_1$&$3-\sqrt5$, strict\\
5&023+01&$\{1,5,6,7,8,9,10,11,12\};\{0,2,3,4\}$&$P_2P_3;T_1$&$0$, strict\\
6&01+02+03&$\{1,5,6,7,8\};\{0,2,3,4,9,10,11,12\}$&$P_2;T_1P_3$&$5-2\sqrt5$, strict\\
7&01+02+13&$\{1,5,6,7,8\};\{0,2,3,4,9,10,11,12\}$&$P_2;T_1P_3$&$5-2\sqrt5$, strict\\
8&01+02+23&$\{5,9,10,11,12\};\{0,1,2,3,4,6,7,8\}$&$P_3;T_0T_1$&$3-\sqrt5$, strict\\
9&01+02+23&$\{6,9,10,11,12\};\{0,1,2,3,4,5,7,8\}$&$P_3;T_0T_1$&$3-\sqrt5$, strict\\
10&023+12&$\{3,4,5\};\{0,1,2,6,7,8,9,10,11,12\}$&$T_1;T_0P_3$&$3-\sqrt5$, strict\\
11&02+03+12&same as type 10&$T_1;T_0P_3$&$3-\sqrt5$, strict\\
12&02+12+23&same as type 10&$T_1;T_0P_3$&$3-\sqrt5$, strict\\
13&02+12+23&same as type 10&$T_1;T_0P_3$&$3-\sqrt5$, strict\\
14&023+12&same as type 10&$T_1;T_0P_3$&$3-\sqrt5$, strict\\
15&02+03+12&same as type 10&$T_1;T_0P_3$&$3-\sqrt5$, strict\\
16&02+12+23&same as type 10&$T_1;T_0P_3$&$3-\sqrt5$, strict\\
17&02+12+23&same as type 10&$T_1;T_0P_3$&$3-\sqrt5$, strict\\
18&02+13+23&$\{3,4,5\};\{0,1,2,6,7,8,9,10,11,12\}$&$T_1;T_0P_2$&$3-\sqrt5$, strict\\
19&02+13+23&same as type 18&$T_1;T_0P_2$&$3-\sqrt5$, strict\\
20&02+13+23&same as type 18&$T_1;T_0P_2$&$3-\sqrt5$, strict
\end{tabular}
\caption{Exact induced-packet certificates for connected shared-cut types
2--20.  The script prints the canonical cycles and independently checks all
vertex subsets.}
\label{tab:3355}
\end{table}

Since $5-2\sqrt5>0$, $3-\sqrt5>0$, and type 5 has zero credit with a strict
triangular packet, \eqref{eq:super} proves $s^+(G)>n$ for types 2--20.  Every
off-core component has a unique core root and is assigned whole to the packet
containing that root.  Thus the finite core certificates apply to arbitrary
attached trees; they are not certificates merely for bare cores.

\subsection{Type 1: the common-root bouquet}

Type 1 has incidence $0123$: all four cycles meet at one cut vertex $o$.  Its
best partition using only Table~\ref{tab:credits} has credit $-1$, so there is
no induced-packet proof of this type.  We give the exact analytic replacement.

\begin{theorem}[Common-root bouquet phase]\label{thm:bouquet}
Suppose $r$ cyclic blocks $C_{\ell_1},\ldots,C_{\ell_r}$ share one common cut
vertex, no $\ell_j$ is divisible by four, and arbitrary trees are attached at
all core vertices.  Then
\begin{equation}\label{eq:bouquet-general}
 s^+(G)\geq n+r-1-
 \sum_{\ell_j\equiv1\ (4)}\left(\sec(\pi/\ell_j)-1\right).
\end{equation}
\end{theorem}

\begin{proof}
Eliminate all trees by \eqref{eq:BP}.  On lobe $j$, list the nonroot
activities as $x_{j,1},\ldots,x_{j,\ell_j-1}$ and define the continuant
\[
 K()=1,\quad K(x_1)=x_1,\quad
 K(x_1,\ldots,x_k)=x_kK(x_1,\ldots,x_{k-1})+K(x_1,\ldots,x_{k-2}).
\]
The exact root-unavailable and root-matched states are
\[
 A_j=K(x_{j,1},\ldots,x_{j,\ell_j-1}),
\]
\[
 B_j=K(x_{j,2},\ldots,x_{j,\ell_j-1})
 +K(x_{j,1},\ldots,x_{j,\ell_j-2}).
\]
If $a$ is the root activity, the matching carrier is
\[
 (\prod_jA_j)\left(a+\sum_jB_j/A_j\right).
\]
At most one whole lobe can be selected in a Sachs subgraph.  Dividing by the
common positive factor $K_{\rm tree}\prod_jA_j$ therefore gives $X+iY$, where
\begin{align*}
 X&=a+\sum_jB_j/A_j+2\sum_{\ell_j\equiv2\ (4)}A_j^{-1}>0,\\
 Y&=2\sum_{\ell_j\equiv1\ (4)}A_j^{-1}
    -2\sum_{\ell_j\equiv3\ (4)}A_j^{-1}.
\end{align*}
For an odd lobe set $c_j=aA_j+B_j$.  Since $X\geq a+B_j/A_j$, positivity and
subadditivity of $\arctan$ give
\[
 -\sum_{\ell_j\equiv3\ (4)}\arctan\frac2{c_j}
 \leq\Theta_G(t)\leq
 \sum_{\ell_j\equiv1\ (4)}\arctan\frac2{c_j}.
\]
All effective activities are at least $t$, so
$c_j\geq Z_{C_{\ell_j}}(t)$.  Hence the upper phase is bounded by the sum of
the corresponding bare $1\pmod4$ cycle phases.  Applying
\eqref{eq:coulson}, using
$D(C_\ell)=-2(\sec(\pi/\ell)-1)$ for $\ell\equiv1\pmod4$, and using
$m=n-1+r$ proves \eqref{eq:bouquet-general}.
\end{proof}

For $C_3,C_3,C_5,C_5$, only the two pentagons are adverse, and
\begin{equation}\label{eq:3355bouquet}
 s^+(G)\geq n+3-2(\sqrt5-2)=n+7-2\sqrt5>n.
\end{equation}
This proves type 1.  It also gives the stronger $C_{333q}$ common-root bound
quoted after Proposition~\ref{prop:333q}.  The positivity of $X$ explains why no cycle of
length $0\pmod4$ is included in this theorem: such a lobe contributes
$-2/A_j$ to $X$.

\begin{proposition}\label{prop:3355}
Every connected cactus with cyclic blocks $C_3,C_3,C_5,C_5$ satisfies
$s^+(G)>n$.
\end{proposition}

\begin{proof}
If the shared-cut graph is disconnected, apply
Proposition~\ref{prop:3355disconnected}.  Otherwise the exact census gives one
of twenty types.  Table~\ref{tab:3355} proves types 2--20, and
\eqref{eq:3355bouquet} proves type 1.  These alternatives include every
connector and every attached tree.
\end{proof}

\section{Exhaustive conclusion}

\begin{theorem}\label{thm:main}
Every connected tetracyclic cactus $G$ on $n$ vertices satisfies
\[
 \boxed{s^+(G)>n}.
\]
\end{theorem}

\begin{proof}
Proposition~\ref{prop:residual} shows that the sharp DNN bound proves the
theorem strictly whenever its excess sum is less than three.  Equality in the
DNN bound can occur only for three triangles and one even cycle: with at most
one triangle the excess sum is at most $1+3\epsilon_5<3$, while with two
triangles the only pair exceeding one is $5,5$ and every other pair is less
than one.  Proposition~\ref{prop:333even} settles this equality case.  The
remaining residual multisets are
$\{3,3,3,q\}$ with $q$ odd and $\{3,3,5,5\}$.  The first complete family is
settled by Proposition~\ref{prop:333q}; the second is settled by
Proposition~\ref{prop:3355}.  The residual proofs exhaust the shared-cut
cluster tree, all common-root and non-common-root incidences, the exact twenty
connected $C_{3355}$ cores, arbitrary bridge paths, and arbitrary rooted trees.
Thus no tetracyclic cactus remains.
\end{proof}

\begin{remark}[What is not used]
The proof never assumes that adding an edge, closing a path into a cycle, or
replacing a tree by a star increases $s^+$.  It uses only induced-subgraph
superadditivity, exact packet theorems, and exact matching-phase identities.
It also does not use a raw coefficientwise phase certificate for all shared
$C_{333q}$ cores: that proposed certificate is false already at $q=5$.
\end{remark}

\appendix
\section{Reproduction and proof objects}\label{app:reproduce}

All paths are relative to the repository root.  Reproduce the twenty-core
$C_{3355}$ census and the exhaustive induced-packet dynamic program with
\begin{verbatim}
python positive-square-energy/experiments/
  c3_c3_c5_c5_induced_packet_partitions.py
\end{verbatim}
The expected final line is
\begin{quote}\ttfamily\small
SUMMARY exact\_positive=18/20 strict\_target=19/20\\
minimum\_score=-1 minimum\_decimal=-1.000000000000
\end{quote}
Here $18/20$ counts literally positive credits; $19/20$ also counts type 5,
whose zero credit is strict.  The missing type is exactly the common-root
bouquet proved by Theorem~\ref{thm:bouquet}.  The script imports
\begin{verbatim}
positive-square-energy/experiments/c3_c3_c5_c5_shared_cluster_certificate.py
\end{verbatim}
which independently generates the twenty canonical core types.  The audit
description and full ledger are
\begin{verbatim}
research/c3-c3-c5-c5-induced-packet-partition-audit-2026-07-26.md
research/c3-c3-c5-c5-tetracyclic-proof-2026-07-26.md
\end{verbatim}

The failed raw $C_{3335}$ coefficient certificate is independently audited by
\begin{verbatim}
python positive-square-energy/experiments/
  c3_c3_c3_c5_shared_cluster_certificate.py
\end{verbatim}
It finds exactly eleven shared core types and only three coefficientwise
positive instances of $\Phi=2R-Z_5I$.  The aggregate exact-ledger SHA-256 is
\begin{verbatim}
f7f39c81aa13443259ca5a644acfeb639041bc44ff393145a3a031520bac4c03
\end{verbatim}
This negative audit is included to make the logical boundary reproducible:
Lemma~\ref{lem:333qterritory} is a territory-credit proof, not a disguised raw
phase assertion.

The analytic common-root calculation is recorded separately at
\begin{verbatim}
research/common-root-cycle-bouquet-phase-2026-07-26.md
\end{verbatim}
The lower-rank packet proofs and sharp DNN proof are
\begin{verbatim}
sharp-cactus-dnn/paper.tex
all-bicyclic-cacti/paper.tex
all-tricyclic-cacti/paper.tex
research/c3-c3-cq-tricyclic-cactus-2026-07-25.md
\end{verbatim}

\begin{thebibliography}{99}
\bibitem{AbiadEtAl}
A.~Abiad, L.~de Lima, D.~N.~Desai, K.~Guo, L.~Hogben, and J.~Madrid,
\emph{Positive and negative square energies of graphs},
Electron. J. Linear Algebra 39 (2023), 307--326.

\bibitem{AKMP}
S.~Akbari, H.~Kumar, B.~Mohar, and S.~Pragada,
\emph{A linear lower bound for the square energy of graphs},
Electron. J. Combin. 32(3) (2025), Paper No. P3.53.

\bibitem{AKMPZ}
S.~Akbari, H.~Kumar, B.~Mohar, S.~Pragada, and S.~Zhang,
\emph{Refinement of a conjecture on positive square energy of graphs},
arXiv:2506.07264, 2025.

\bibitem{LTZ}
Y.~Liu, Q.~Tang, and S.~Zhang,
\emph{The positive and negative square-energy conjecture},
arXiv:2607.18031, 2026.

\bibitem{NingZeng}
B.~Ning and J.~Zeng,
\emph{A proof of a conjecture on positive and negative square energies of
unicyclic graphs}, arXiv:2605.24668, 2026.

\bibitem{SharpDNN}
Companion manuscript, \emph{The Optimized Liu--Tang--Zhang Constant for
Cactus Graphs}, 2026; repository path \texttt{sharp-cactus-dnn/paper.tex}.

\bibitem{AllBicyclic}
Companion manuscript, \emph{Positive Square Energy of Every Bicyclic Cactus},
2026; repository path \texttt{all-bicyclic-cacti/paper.tex}.

\bibitem{AllTricyclic}
Companion manuscript, \emph{Positive Square Energy of Every Tricyclic Cactus},
2026; repository path \texttt{all-tricyclic-cacti/paper.tex}.
\end{thebibliography}

\end{document}
